\documentclass[letterpaper]{article}
\usepackage[submission]{aaai2027}
\usepackage[hyphens]{url}
\urlstyle{rm}
\def\UrlFont{\rm}
\usepackage{natbib}
\usepackage{booktabs}
\frenchspacing

\pdfinfo{
/TemplateVersion (2027.1)
}

\setcounter{secnumdepth}{0}
\newcommand{\PublicSkillCount}{79,141}
\newcommand{\PublicTaskCount}{73}
\newcommand{\PublicVariantCount}{219}
\newcommand{\PublicGoldCount}{182}
\newcommand{\WordEnHit}{37.0\%}
\newcommand{\WordEnRecall}{46.8\%}
\newcommand{\WordEnNdcg}{0.370}
\newcommand{\WordZhHit}{8.2\%}
\newcommand{\WordZhRecall}{17.7\%}
\newcommand{\WordZhNdcg}{0.116}
\newcommand{\WordInstHit}{8.2\%}
\newcommand{\WordInstRecall}{17.7\%}
\newcommand{\WordInstNdcg}{0.114}
\newcommand{\WordLatency}{5.7}
\newcommand{\BmEnHit}{46.6\%}
\newcommand{\BmEnRecall}{46.0\%}
\newcommand{\BmEnNdcg}{0.391}
\newcommand{\BmZhHit}{19.2\%}
\newcommand{\BmZhRecall}{21.7\%}
\newcommand{\BmZhNdcg}{0.168}
\newcommand{\BmInstHit}{12.3\%}
\newcommand{\BmInstRecall}{17.1\%}
\newcommand{\BmInstNdcg}{0.127}
\newcommand{\BmLatency}{5.6}
\newcommand{\CharEnHit}{49.3\%}
\newcommand{\CharEnRecall}{50.1\%}
\newcommand{\CharEnNdcg}{0.432}
\newcommand{\CharZhHit}{21.9\%}
\newcommand{\CharZhRecall}{31.4\%}
\newcommand{\CharZhNdcg}{0.235}
\newcommand{\CharInstHit}{24.7\%}
\newcommand{\CharInstRecall}{29.4\%}
\newcommand{\CharInstNdcg}{0.236}
\newcommand{\CharLatency}{14.8}
\newcommand{\EfiveEnHit}{13.7\%}
\newcommand{\EfiveEnRecall}{15.5\%}
\newcommand{\EfiveEnNdcg}{0.119}
\newcommand{\EfiveZhHit}{0.0\%}
\newcommand{\EfiveZhRecall}{0.0\%}
\newcommand{\EfiveZhNdcg}{0.000}
\newcommand{\EfiveInstHit}{0.0\%}
\newcommand{\EfiveInstRecall}{0.0\%}
\newcommand{\EfiveInstNdcg}{0.000}
\newcommand{\EfiveLatency}{12.0}
\newcommand{\SrEnHit}{45.2\%}
\newcommand{\SrEnRecall}{50.5\%}
\newcommand{\SrEnNdcg}{0.412}
\newcommand{\SrZhHit}{41.1\%}
\newcommand{\SrZhRecall}{43.8\%}
\newcommand{\SrZhNdcg}{0.366}
\newcommand{\SrInstHit}{26.0\%}
\newcommand{\SrInstRecall}{38.9\%}
\newcommand{\SrInstNdcg}{0.298}
\newcommand{\SrLatency}{140.8}
\newcommand{\FusionEnHit}{49.3\%}
\newcommand{\FusionEnRecall}{54.9\%}
\newcommand{\FusionEnNdcg}{0.456}
\newcommand{\FusionZhHit}{42.5\%}
\newcommand{\FusionZhRecall}{46.0\%}
\newcommand{\FusionZhNdcg}{0.375}
\newcommand{\FusionInstHit}{39.7\%}
\newcommand{\FusionInstRecall}{42.1\%}
\newcommand{\FusionInstNdcg}{0.353}
\newcommand{\FusionLatency}{155.6}
\newcommand{\SparseLatencyRange}{5.6--14.8}
\newcommand{\SrVsCharZhSignificance}{$p=0.013$, exact McNemar, 21 gains vs. 7 losses}
\newcommand{\SrVsCharInstSignificance}{$p=1.000$, exact McNemar, 14 gains vs. 13 losses}


\title{Technical Supplement: Bilingual Skill Routing at 80K Scale}
\author{Anonymous Submission}
\affiliations{}

\begin{document}
\maketitle

\section{Public Benchmark Construction}

The SkillRouter Eval Core release \cite{zheng2026skillrouter} contains 87 tasks
and two candidate pools.
We use its 79,141-record hard tier and omit two source-excluded and 12
\texttt{generic\_only} tasks, leaving \PublicTaskCount{} tasks and
\PublicGoldCount{} distinct core gold skill identifiers.  Every gold identifier
is present in the candidate pool.

The released multilingual file stores the source fields and three query
columns: exact English, Simplified Chinese, and the same Chinese task with one
fixed bilingual routing instruction.
The first translation pass was generated by one isolated language-model run.
A second isolated run compared every translated row with the exact English
source and corrected all 73 rows.  The audit focused on omitted steps,
negations, numeric constraints, paths, code symbols, fenced code blocks, and
output schemas.  A deterministic validator then checks:

exact English preservation, complete task coverage without duplicates,
Chinese-script presence in translated queries, both Chinese and Latin scripts
in instruction-prefixed queries, exact prefix equality, and the absence of
email-shaped personal identifiers.

The data card and main paper disclose that these are not professional human
translations.

\section{Candidate Representation}

For all methods, the candidate text contains only the skill name and
description.  The public records also provide a body, but using it would change
the research question from pre-activation metadata routing to full-text
routing.  Word and character methods join fields with a period.  Generic
multilingual-E5 receives \texttt{passage: name. description}.  The
SkillRouter-domain encoder receives \texttt{name | description}, matching the
field delimiter in its public model card while omitting the body.

\section{Hyperparameters}

Word TF--IDF uses 1--2 grams, sublinear term frequency, and a 300K feature
cap.  BM25 uses unigrams, $k_1=1.2$, $b=0.75$, and a 300K cap.  Character
TF--IDF uses 2--5 grams, minimum document frequency 2, and a 300K cap.
mE5-small uses 384 dimensions, length 512, and batch 128; SR-Emb uses 1024
dimensions, length 512, and batch 8.  RRF combines the top 20 SR and character
results with constant 60.  Bootstrap confidence intervals use 10,000 samples
and seed 20270827.  No task-specific model training is performed.
Both Hugging Face models are pinned to exact 40-character revisions in
\texttt{config/model\_revisions.json}; cache metadata additionally binds the
model revision, ordered skill identifiers, formatted document text, matrix
shape, and output checksum.
Exactly one value was tried for each public-benchmark setting; values were
fixed before evaluation from standard formulations, model cards, and the
source benchmark protocol rather than selected on test performance.

The public dataset and model caches are downloaded or generated once.  Query
latency includes query vectorization or model encoding plus exact search over
the cached candidate matrix.  Sparse index fitting and dense cache
load-and-validation time are recorded as \texttt{setup\_seconds}.  Offline
candidate embedding generation is a separate one-time step and is not timed in
the reported result.  Equal scores use ascending skill identifier as a stable
secondary key.
The public SkillRouter dataset is pinned to an immutable Hugging Face commit in
\texttt{config/dataset\_revisions.json}, and every downloaded file is checked
against a committed SHA-256 digest.

\section{Reproduction}

The release includes:
\begin{itemize}
  \item SHA-256 checksums and a downloader for the public SkillRouter hard tier;
  \item an audited multilingual JSONL file and deterministic validator;
  \item exact dependency locks;
  \item scripts for both embedding caches and all retrieval methods;
  \item result JSON, rankings, figures, and task-level error analysis; and
  \item a Makefile that compiles the paper, checklist, supplement, and archives.
\end{itemize}

The full workflow is:
\begin{verbatim}
make setup
make data validate-data
make cache-e5 cache-skillrouter
make experiments test
make submission
\end{verbatim}

The reported run used a 14-core Apple M4 Pro with 64\,GB memory, Apple MPS,
macOS 26.6.1, Python 3.12.14, PyTorch 2.13.0, scikit-learn 1.9.0, and
sentence-transformers 5.7.0.  Exact versions are stored in
\texttt{requirements-aaai-lock.txt}.

\section{License and Attribution}

The source task text is derived from SkillsBench \cite{li2026skillsbench},
whose repository uses Apache License 2.0, and is redistributed by SkillRouter
Eval Core.  The multilingual
variants are marked as modifications and retain upstream attribution.  The
large candidate pool and model caches are not committed; the release downloads
them from their official locations.  Code written for this study uses the MIT
License.  Third-party skill descriptions retain their original terms.

\bibliography{references}
\end{document}
